\SourceLeaf{29}
\fancyhead[L]{\small ALLGEMEINES GRUNDGESETZ DER MECHANIK.}
\pdfbookmark[0]{Über ein neues allgemeines Grundgesetz der Mechanik}{mechanik}
\begin{center}
ÜBER EIN NEUES\par
\bigskip
ALLGEMEINES GRUNDGESETZ DER MECHANIK.\par
\medskip
\rule{.16\textwidth}{.3pt}
\end{center}
\expandafter\def\csname HistoricalNoteText2\endcsname{Ich möchte hier jedoch bemerken, dass ich die Art, wie ein anderer großer Geometer versucht hat, \textsc{Huyghens}’ Gesetz für die außerordentliche Brechung des Lichts in doppelbrechenden Kristallen mithilfe des Prinzips der kleinsten Wirkung zu beweisen, nicht befriedigend finde. Ob dieses Prinzip anwendbar ist, hängt in der Tat wesentlich vom Prinzip der Erhaltung der lebendigen Kräfte ab. Danach sind die Geschwindigkeitsbeträge der bewegten materiellen Punkte allein durch ihre Orte bestimmt, ohne dass die Bewegungsrichtung darauf Einfluss haben kann. Ein solcher Einfluss wird aber in dem erwähnten Versuch vorausgesetzt. Mir scheint, dass im Rahmen der Emissionstheorie alle Bemühungen, die Erscheinungen der Doppelbrechung mit den allgemeinen Gesetzen der Dynamik zu verbinden, so lange erfolglos bleiben müssen, wie man die Lichtteilchen lediglich als Punkte betrachtet.}
\par
Bekanntlich macht das Prinzip der virtuellen Geschwindigkeiten die gesamte Statik zu einer mathematischen Aufgabe, und \textsc{d’Alembert}s Prinzip für die Dynamik führt diese wiederum auf die Statik zurück. Es liegt daher in der Natur der Sache, dass es kein neues Grundprinzip für die Lehre von Bewegung und Gleichgewicht geben \textit{kann}, das seinem Inhalt nach nicht schon in diesen beiden enthalten und aus ihnen ableitbar wäre. Trotzdem scheint dieser Umstand noch nicht jedes neue Prinzip wertlos zu machen. Es wird stets interessant und lehrreich bleiben, den Naturgesetzen einen neuen, günstigen Blickwinkel abzugewinnen, sei es, weil sich aus diesem Blickwinkel die eine oder andere Aufgabe leichter lösen lässt, oder weil sich eine besondere Angemessenheit zeigt. Der große Geometer, der das Gebäude der Mechanik so glänzend auf dem Fundament des Prinzips der virtuellen Geschwindigkeiten errichtet hat, hat es nicht verschmäht, \textsc{Maupertuis}’ Prinzip der kleinsten Wirkung präziser und allgemeiner zu fassen. Dieses Prinzip lässt sich bisweilen mit großem Vorteil anwenden\HistoricalNote{2}.
